% !TeX root = ../nc-quadric-avatar.tex
% 2026-11-22 / smooth ruling interchange through a four-sheeted half-period torsor
\section{滑らかなnoncubicのruling交換と半周期torsor}
\label{sec:qa-half-period-ruling}

前節は一般係数体でruling交換が頂点固定のMorita
自己同値にならないことを示した。
一方、幾何的に四項を交換するための
\emph{有限\'{e}tale被覆}は、Jacobianの半周期から構成できる。
ここで点曲線が滑らかな開に限定し、
nodal/$I_2$への同じ次数四の延長は主張しない。

\begin{theorem}[半周期による二pencil交換の次数四torsor]
\label{thm:qa-half-period-ruling-torsor}
$U^\circ\subset U_2$を、二つの二次形式の四根が相異なる開とし、
$E\to U^\circ$を滑らかな種数一の点曲線、
$J=\Pic^0_{E/U^\circ}$をその相対Jacobianとする。
前節のBrauer gerbe上で定義した
$P=[GF^{-1}]$は、根の選択や線束の
基底捩れによらない非零な
$P\in J(U^\circ)[2]$へ降下する。
$[2]:J\to J$のこのsection上のfiber
\begin{equation}
 T_P=U^\circ\times_{P,J,[2]}J
 \xrightarrow{\quad}U^\circ
 \label{eq:qa-half-period-four-cover}
\end{equation}
は有限\'{e}tale次数四で、
$J[2]$のtorsorである。

$T_P$上のtautological sectionを$t$とする。
種数一torsor $E_{T_P}$に対する
Jacobianの平行移動$\tau_t$は、
相対Picard類について
\begin{equation}
 [\tau_t^*F]=[G],\quad
 [\tau_t^*G]=[F],\quad
 [\tau_t^*H]=[H],\quad
 [\tau_t^*(HG)]=[HF]
 \label{eq:qa-half-period-exchange}
\end{equation}
を満たす。
従って相対四項
$(\OO_E,F,H,HG)$と$(\OO_E,G,H,HF)$は、
$T_P$をBrauer gerbeへ引き戻した後、
必要な基底線束の捩れを除けば
$\tau_t^*$によって全Homと積を保って対応する。
実際の各頂点線束の同型を同時に選ぶには、
さらにfppf局所化してよい。
これは幾何的finite Endの比較であって、
相対AS mutationや相対pointed$\qgr$同値を証明しない。

一般係数体$K_{12}$上では$T_P$に有理点がない。
従って上記の半周期による四頂点交換を
一般係数体上で一点の選択として実現することはできない。
\end{theorem}
\begin{proof}
$E$は元の二conicの係数空間$U^\circ$上で定義される。
gerbe上の$F,G$の相対Picard類は
線束の基底からの捩れに不変であり、
各$\Gm$慣性も相対Picard類には自明に作用する。
よって$P=[GF^{-1}]$は$J$のsectionへ降下する。
各幾何fiberで非零二捩れであることは
補題~\ref{lem:qa-pointed-all-rootpair-geometry}による。
標数零の楕円曲線の$[2]$は有限\'{e}tale次数四であるため、
base changeした$T_P$も同じ性質を持ち、
各fiberは$J[2]$の自由かつ推移的なtorsorである。

滑らかな種数一曲線$E_s$を幾何点$s$上でとり、
一時的に原点を選ぶ。
次数$d$の線束$L$は次数$d$の因子の線型同値類であり、
平行移動によるpullbackは
\[
 [\tau_t^*L]-[L]=-d t\quad\text{in }\Pic^0(E_s)
 \tag{Translation-degree}\label{eq:qa-half-period-degree}
\]
となる。実際、$\sum_\nu x_\nu$のpullbackは
$\sum_\nu(x_\nu-t)$で、両者の差は$-dt$である。
この公式は原点の選択によらないため相対Picardに降下する。
$F,G$の次数はともに二、$H$の次数は四であり、
$2t=P$かつ$2P=0$だから
\[
 -2t=-P=P,\qquad -4t=-2P=0.
\]
従って$\tau_t^*$は$F$を$FP=G$に、
$G$を$GP=F$に、$H$を$H$に送る。
積についても同じであり、
式~\eqref{eq:qa-half-period-exchange}を得る。
曲線の自己同型のpullbackは射の合成を保つので、
各頂点の線束同型をfppf局所で選べば
有限Endの全合成が比較できる。
ただしその選択のdescent datumまではここで作っていない。

最後に$K_{12}$上で$2t=P$となる
$K_{12}$-有理点$t$が存在すると仮定する。
この$t$による平行移動は
$K_{12}$-自己同型$\tau_t$である。
幾何Picard類$F,G$のHochschild--Serre境界は
前節の命題~
\ref{prop:qa-universal-conic-brauer-obstruction}により
$\partial_E[F]=\alpha_c$、
$\partial_E[G]=\alpha_d$である。
この境界は基礎体上の自己同型pullbackと可換し、
右端の$\operatorname{Br}(K_{12})$には自明に作用する。
$\tau_t^*[F]=[G]$から
$\alpha_c=\alpha_d$が従うが、
同命題で$\alpha_c-\alpha_d\ne0$と証明している。
矛盾するので$T_P(K_{12})=\varnothing$である。
\end{proof}

\begin{remark}[異なる次数四被覆と特異境界]
\label{rem:qa-half-period-versus-root-cover}
半周期torsor $T_P$は$J[2]$-torsorであり、
前節の「各二次形式から分岐根を一つずつ選ぶ」
有限\'{e}tale次数四被覆とは定義が異なる。
両被覆の同型、OUの一般次数四射との比較、
およびそのGaloisモノドロミーの一致は未確認である。
nodal/$I_2$では相対$\Pic^0$の$[2]$の幾何次数が
smooth fiberの四から変わるため、
上記の同じ次数四の\'{e}tale torsorを
境界まで自動的に延長してはならない。
\end{remark}
